\documentclass[11pt,a4paper]{article}
% Rapport du mini projet 16. Compiler avec : xelatex rapport.tex (deux fois)
% Les passages À REMPLIR sont surlignés en jaune : à compléter avec vos résultats.
\usepackage{fontspec}
\usepackage{amsmath,amssymb}
\usepackage{babel}
\babelprovide[import,main]{french}
\usepackage[a4paper,margin=2.3cm]{geometry}
\usepackage{xcolor,graphicx,booktabs,tabularx,array,float,caption,enumitem}
\usepackage[most]{tcolorbox}
\usepackage[hidelinks]{hyperref}
\setlength{\parskip}{0.4em}
\sloppy\setlength{\emergencystretch}{3em}
\newcolumntype{L}{>{\raggedright\arraybackslash}X}
\AtBeginDocument{\renewcommand{\tablename}{Tableau}}

% Marqueurs de ce qu'il reste à remplir
\definecolor{jaune}{HTML}{FFF2A8}
\definecolor{bordjaune}{HTML}{C98500}
\newcommand{\ar}[1]{\colorbox{jaune}{\textbf{\small [#1]}}}            % court, dans une phrase
\newtcolorbox{aremplir}{colback=jaune,colframe=bordjaune,boxrule=0.5pt,arc=1pt,
  left=5pt,right=5pt,top=3pt,bottom=3pt,fontupper=\small,
  title={À REMPLIR},fonttitle=\small\bfseries,coltitle=black,colbacktitle=jaune!70!bordjaune}
\newcommand{\redaction}[1]{\hfill{\small\color{gray}(rédaction : #1)}}   % à retirer avant le rendu
% Figure : insérée si le fichier existe dans resultats/, sinon un cadre indique quoi produire
\newcommand{\figureres}[3]{\begin{figure}[H]\centering
  \IfFileExists{resultats/#1}{\includegraphics[width=#2\linewidth]{resultats/#1}}%
  {\fcolorbox{bordjaune}{jaune}{\parbox{0.8\linewidth}{\centering\small\vspace{1.5em}
   \textbf{FIGURE À INSÉRER} : \texttt{resultats/\detokenize{#1}}\\ (produite par le script indiqué dans la légende)\vspace{1.5em}}}}
  \caption{#3}\end{figure}}

\begin{document}

%=====================================================================
\begin{titlepage}
\centering
{\large École Polytechnique de Thiès\\Département Génie Informatique et Télécommunications\par}
\vspace{1cm}
{\large DIC1 GIT, 2026-2027 · Traitement du signal\par}
\vspace{3cm}
{\LARGE\bfseries Reconnaissance automatique de mots isolés\\ dans un signal bruité et variable en locuteur\par}
\vspace{0.8cm}
{\large Mini projet 16 : volet classique sous MATLAB et benchmark des méthodes modernes\par}
\vspace{3cm}
{\large Réalisé par :\par}
\vspace{0.3cm}
\ar{Prénom NOM (personne A)}\\[0.3em]
\ar{Prénom NOM (personne B)}\\[0.3em]
\ar{Prénom NOM (personne C)}\par
\vspace{2cm}
{\large Sous la direction de : Pr N.F. NGOM\par}
\vfill
{\large 16 novembre 2026\par}
\end{titlepage}

%=====================================================================
\section*{Résumé}
\begin{aremplir}
À écrire en dernier (8 à 10 lignes) : la question de recherche, les méthodes comparées (corrélation, DTW, HMM discret, HMM-GMM sur MFCC), la base (Speech Commands v2, 10 mots, test indépendant du locuteur), les deux ou trois chiffres principaux (meilleure accuracy propre, accuracy à 0 dB, écart dépendant / indépendant), et la conclusion du benchmark.
\end{aremplir}

\tableofcontents
\newpage

%=====================================================================
\section{Introduction}\redaction{tous}

La reconnaissance automatique de la parole permet de commander une machine par la voix : assistants vocaux, domotique, systèmes embarqués, accessibilité. Dans sa forme la plus simple, la reconnaissance de \emph{mots isolés}, le système doit identifier une commande parmi un petit vocabulaire. Même dans ce cadre restreint, deux difficultés majeures subsistent : le \textbf{bruit} de l'environnement, qui masque une partie du signal, et la \textbf{variabilité entre locuteurs} (hauteur de voix, accent, débit, micro), qui fait qu'un même mot ne produit jamais deux fois le même signal.

Ce mini projet répond à la question de recherche suivante :
\begin{center}
\emph{Comment reconnaître automatiquement des mots ou des commandes vocales à partir d'un signal audio bruité et variable en locuteur ?}
\end{center}

Pour y répondre, nous avons implémenté sous MATLAB une chaîne de reconnaissance classique (prétraitement, extraction de MFCC, reconnaissance par corrélation, par alignement temporel dynamique et par modèles de Markov cachés), que nous avons évaluée sur la base Google Speech Commands selon un protocole indépendant du locuteur, en conditions propres et bruitées. Nous avons ensuite comparé ces résultats à ceux publiés pour les méthodes modernes fondées sur l'apprentissage profond.

Le rapport est organisé selon les trois axes demandés : la section~\ref{sec:etat} situe les méthodes ; la section~\ref{sec:methodes} présente la chaîne classique et ses hypothèses (axe 1) ; la section~\ref{sec:protocole} décrit le protocole ; la section~\ref{sec:resultats} présente les performances (axe 2) ; la section~\ref{sec:benchmark} compare nos résultats à l'état de l'art ; la section~\ref{sec:perspectives} propose des perspectives d'applications (axe 3).

%=====================================================================
\section{État de l'art}\label{sec:etat}\redaction{C}

\paragraph{Les méthodes classiques.} Les premiers systèmes de reconnaissance de mots isolés (années 1970) comparaient le mot prononcé à des enregistrements de référence par \textbf{alignement temporel dynamique} (DTW, Sakoe et Chiba, 1978), qui compense les différences de débit entre deux prononciations. La représentation du signal par les \textbf{coefficients cepstraux sur l'échelle Mel} (MFCC, Davis et Mermelstein, 1980) s'est imposée car elle décrit l'enveloppe spectrale, liée au son prononcé, en écartant la fréquence fondamentale, liée au locuteur. Hermansky (1990) a proposé les coefficients \textbf{PLP}, qui intègrent davantage de propriétés de l'audition. À partir des années 1980, les \textbf{modèles de Markov cachés} à émissions par mélanges de gaussiennes (HMM-GMM, Rabiner, 1989) ont remplacé la DTW : au lieu de mémoriser des exemples, ils apprennent un modèle statistique de chaque mot ou de chaque phonème.

\paragraph{Les méthodes modernes.} Vers 2012, les réseaux de neurones profonds ont remplacé les GMM dans les systèmes hybrides \textbf{DNN-HMM} (Hinton et al., 2012), puis les \textbf{TDNN} de la boîte à outils Kaldi (Peddinti et al., 2015). Les modèles de bout en bout comme le \textbf{Conformer} (Gulati et al., 2020) ont ensuite supprimé le HMM. Les modèles \textbf{auto-supervisés} wav2vec 2.0 (Baevski et al., 2020) et HuBERT (Hsu et al., 2021) apprennent à partir de grandes quantités d'audio non transcrit, tandis que \textbf{Whisper} (Radford et al., 2023) est entraîné sur 680\,000 heures d'audio transcrit. Enfin, la détection de mots-clés embarquée (TinyML) vise des modèles de quelques milliers de paramètres pour microcontrôleurs (Zhang et al., 2017 ; Kim et al., 2021).

%=====================================================================
\section{Méthodes classiques et hypothèses (axe 1)}\label{sec:methodes}

La figure~\ref{fig:chaine} résume la chaîne implémentée. Tous les réglages sont regroupés dans un fichier unique (\texttt{configProjet.m}) ; le tableau~\ref{tab:reglages} les récapitule.

\begin{figure}[H]\centering
\fbox{\parbox{0.92\linewidth}{\centering\small
Signal (16 kHz) $\rightarrow$ centrage, normalisation $\rightarrow$ VAD $\rightarrow$ préaccentuation $\rightarrow$ trames de 25 ms (pas 10 ms), Hamming $\rightarrow$ FFT 512 points $\rightarrow$ 26 filtres Mel $\rightarrow$ log $\rightarrow$ DCT (13 MFCC) $\rightarrow$ deltas, delta-deltas, CMVN (39 valeurs par trame) $\rightarrow$ corrélation / DTW / HMM $\rightarrow$ mot reconnu}}
\caption{Chaîne de reconnaissance implémentée sous MATLAB.}\label{fig:chaine}
\end{figure}

\begin{table}[H]\centering\small
\begin{tabularx}{\linewidth}{llLL}
\toprule
Étape & Paramètre & Valeur & Justification \\
\midrule
Échantillonnage & $f_e$ & 16 kHz & théorème de Shannon, bande utile jusqu'à 8 kHz (fricatives) \\
Préaccentuation & $\alpha$ & 0,97 & compense la pente spectrale d'environ $-6$ dB/octave de la voix \\
Trames & durée, pas & 25 ms, 10 ms & quasi-stationnarité de la parole sur 20 à 30 ms \\
Fenêtre & type & Hamming & premier lobe secondaire à $-43$ dB, limite la fuite spectrale \\
Banc Mel & filtres, FFT & 26, 512 points & résolution fréquentielle proche de celle de l'oreille \\
MFCC & coefficients & 13 (+ 13 $\Delta$ + 13 $\Delta\Delta$) & enveloppe spectrale ; dynamique temporelle \\
VAD & seuil & bruit de fond + 12 dB & \ar{à ajuster} \\
DTW & décision & 1-NN, distance normalisée par $I+J$ & \ar{bande retenue} \\
HMM-GMM & états, gaussiennes & \ar{N} états, \ar{K} gaussiennes & choisis sur l'ensemble de validation (section~\ref{sec:resultats}) \\
\bottomrule
\end{tabularx}
\caption{Réglages de la chaîne de traitement.}\label{tab:reglages}
\end{table}

\subsection{Prétraitement}\redaction{A}

La parole n'est pas stationnaire, mais les articulateurs évoluent lentement : sur 20 à 30 ms, le conduit vocal varie peu. L'\textbf{hypothèse de quasi-stationnarité à court terme} justifie le découpage en trames de 25 ms, décalées de 10 ms, sur lesquelles on applique la transformée de Fourier. Chaque trame est pondérée par une fenêtre de Hamming pour limiter la fuite spectrale. La \textbf{préaccentuation} $y(n) = x(n) - 0{,}97\,x(n-1)$ est un filtre RIF du premier ordre dont le gain vaut $-30{,}5$ dB à 0 Hz et $+5{,}9$ dB à 8 kHz : elle rééquilibre le spectre en relevant les formants aigus. La \textbf{détection d'activité vocale} (VAD) compare l'énergie de chaque trame au niveau du bruit de fond, estimé par le 10\textsuperscript{e} percentile des énergies du clip ; les détections de moins de 50 ms sont supprimées et une marge de 20 ms est ajoutée de part et d'autre du mot.

\begin{aremplir}
Comparaison de votre VAD avec \texttt{detectSpeech} (script \texttt{verifierVAD.m}) : écart moyen des bornes au début et à la fin, et un ou deux cas d'échec commentés (fricative ratée, bruit fort).
\end{aremplir}

\subsection{Extraction des caractéristiques}\redaction{A}

Le spectre de puissance de chaque trame est intégré dans 26 filtres triangulaires régulièrement espacés sur l'échelle Mel, $m = 2595 \log_{10}(1 + f/700)$, qui reproduit la résolution décroissante de l'oreille en haute fréquence. Le logarithme des énergies transforme le produit source-filtre du modèle de production de la parole en une somme ; la transformée en cosinus discrète (DCT) décorrèle ensuite ces log-énergies, et la conservation des 13 premiers coefficients ne garde que l'enveloppe spectrale. C'est le principe du \textbf{cepstre} : les premiers coefficients décrivent le conduit vocal (le son prononcé), tandis que la périodicité de la source (la hauteur de la voix, propre au locuteur) apparaît à des quéfrences plus élevées et est éliminée. Les \textbf{coefficients dynamiques} (deltas et delta-deltas, calculés par régression sur 5 trames) ajoutent l'information de transition entre sons. Enfin, la \textbf{normalisation CMVN} centre et réduit chaque coefficient sur l'enregistrement, ce qui supprime l'effet d'un filtre de canal constant (micro, pièce), très variable dans une base enregistrée par plus de 2\,600 personnes.

\begin{aremplir}
Validation de votre implémentation (script \texttt{verifierMFCC.m}) : comparaison de votre banc de filtres avec \texttt{designAuditoryFilterBank} et corrélation de vos MFCC avec ceux de \texttt{mfcc} (valeurs obtenues). Si les PLP ont été implémentés, les décrire ici en 3 ou 4 lignes ; sinon, les présenter brièvement en théorie.
\end{aremplir}

\subsection{Reconnaissance par comparaison : corrélation et DTW}\redaction{B}

Dans ces deux méthodes, le système conserve 20 enregistrements de référence par mot, tous prononcés par des locuteurs différents, et attribue au mot inconnu le mot de la référence la plus proche (plus proche voisin).

Le \textbf{template matching par corrélation} rééchantillonne chaque suite de MFCC à une longueur fixe de 50 trames par interpolation linéaire, puis calcule le coefficient de corrélation normalisé. Il suppose implicitement qu'un changement de débit étire uniformément tous les sons du mot.

La \textbf{DTW} lève cette hypothèse : elle cherche le chemin d'alignement monotone entre les trames des deux mots qui minimise la somme des distances euclidiennes locales, par programmation dynamique :
\[
D(i, j) = d(\mathbf{a}_i, \mathbf{b}_j) + \min\{D(i-1, j),\ D(i-1, j-1),\ D(i, j-1)\}.
\]
La distance finale est normalisée par $I + J$ pour ne pas favoriser les références courtes. Ses hypothèses sont l'alignement monotone des débuts et des fins de mots (d'où l'importance de la VAD) et l'égale importance de toutes les trames.

\begin{aremplir}
Validation de votre DTW (script \texttt{testDTWjouet.m}) : accord avec les exemples calculés à la main et avec la fonction \texttt{dtw} de MATLAB. Préciser si la bande de Sakoe-Chiba a été utilisée et pourquoi.
\end{aremplir}

\subsection{Modèles de Markov cachés}\redaction{B}

Chaque mot est modélisé par un HMM gauche-droite de $N$ états : on ne peut que rester dans un état ou passer au suivant, ce qui traduit l'ordre fixe des sons du mot. Chaque état émet les vecteurs de caractéristiques selon une loi de probabilité, et un mot inconnu $O$ est attribué au mot $w$ qui maximise $\log P(O \mid \lambda_w)$ (règle de Bayes avec des mots équiprobables).

Nous avons implémenté deux variantes. Le \textbf{HMM discret} remplace chaque trame par l'indice du plus proche de 64 vecteurs types obtenus par k-moyennes (quantification vectorielle) ; il est appris par l'algorithme de Baum-Welch (\texttt{hmmtrain}) et évalué par l'algorithme Forward (\texttt{hmmdecode}). Le \textbf{HMM-GMM} modélise directement les vecteurs continus par un mélange de gaussiennes à covariances diagonales dans chaque état ; il est appris par \textbf{entraînement de Viterbi} : découpage initial uniforme, apprentissage d'un GMM par état (\texttt{fitgmdist}), réalignement des trames par l'algorithme de Viterbi, et itération. Cette méthode est une approximation de Baum-Welch, plus simple à mettre en œuvre.

Les hypothèses du HMM sont la propriété de Markov d'ordre 1, l'indépendance des trames conditionnellement à l'état (en partie compensée par les deltas), la modélisation des émissions par des gaussiennes et une durée géométrique dans chaque état. Les HMM à états phonétiques, qui modélisent chaque phonème par un HMM à trois états et construisent les mots par concaténation, constituent la généralisation naturelle de cette approche aux grands vocabulaires \ar{implémentés ou non : préciser}.

%=====================================================================
\section{Protocole expérimental}\label{sec:protocole}\redaction{C}

\paragraph{Données.} Nous utilisons la version 2 de Google Speech Commands (Warden, 2018) : 105\,829 enregistrements d'une seconde, échantillonnés à 16 kHz, de 35 mots prononcés par 2\,618 locuteurs avec leur propre matériel. Nous retenons les 10 mots du protocole standard : \emph{yes, no, up, down, left, right, on, off, stop, go}. Le partage officiel, calculé par hachage sur l'identifiant du locuteur, garantit qu'aucun locuteur n'apparaît dans deux ensembles ; nous l'avons vérifié sur les fichiers (tableau~\ref{tab:donnees}).

\begin{table}[H]\centering\small
\begin{tabular}{lrrr}
\toprule
 & Apprentissage & Validation & Test \\
\midrule
Fichiers, 10 mots & 30\,769 & 3\,703 & 4\,074 \\
Fichiers, base entière (35 mots) & 84\,843 & 9\,981 & 11\,005 \\
Locuteurs, base entière & 2\,112 & 256 & 250 \\
\bottomrule
\end{tabular}
\caption{Partage officiel de Speech Commands v2. Aucun locuteur n'est commun à l'apprentissage et au test.}\label{tab:donnees}
\end{table}

\paragraph{Sous-ensembles.} Pour limiter le temps de calcul de la DTW, nous tirons avec une graine fixe (reproductibilité) : 100 fichiers de test par mot (1\,000 au total), 30 fichiers de validation par mot, 20 références par mot de locuteurs tous différents, et 200 exemples d'apprentissage par mot pour les HMM. \textbf{Toutes les méthodes sont évaluées sur les mêmes 1\,000 fichiers de test.} Un test \emph{dépendant du locuteur} est construit avec les autres répétitions des locuteurs de référence (\ar{nombre} fichiers).

\paragraph{Bruit.} Les conditions bruitées sont obtenues en ajoutant au clip complet, avant toute la chaîne de traitement, un bruit blanc ou un bruit réel (robinet) tiré du dossier \texttt{\_background\_noise\_}, à un SNR de 30, 20, 10, 5 et 0 dB. Le gain du bruit est $g = \sqrt{P_s / (P_b\,10^{\text{SNR}/10})}$, où la puissance du signal $P_s$ est calculée sur tout le clip, silences compris : le SNR effectif sur la partie parlée est donc légèrement supérieur au SNR annoncé. Les références restent propres (conditions non appariées).

\paragraph{Métriques.} Accuracy et intervalle de confiance à 95 \% ($\pm 1{,}96\sqrt{p(1-p)/n}$), matrice de confusion, précision, rappel et F1 par mot, temps de calcul mesuré par \texttt{timeit}, facteur temps réel et mémoire. Pour des mots isolés, chaque erreur est une substitution : le WER vaut $1 - \text{accuracy}$. Machine utilisée : \ar{processeur, mémoire, version de MATLAB}.

%=====================================================================
\section{Résultats et analyse des performances (axe 2)}\label{sec:resultats}\redaction{C, avec A et B}

\subsection{Performances en conditions propres}

\begin{table}[H]\centering\small
\begin{tabularx}{\linewidth}{lLLLL}
\toprule
Méthode & Accuracy (IC 95 \%) & F1 macro & Temps / mot & Mémoire ou paramètres \\
\midrule
Corrélation + MFCC & \ar{\%} & \ar{ } & \ar{ms} & \ar{ } \\
DTW + MFCC & \ar{\%} & \ar{ } & \ar{ms} & \ar{ } \\
DTW + MFCC + $\Delta$ + CMVN & \ar{\%} & \ar{ } & \ar{ms} & \ar{ } \\
HMM discret (64 symboles) & \ar{\%} & \ar{ } & \ar{ms} & \ar{ } \\
HMM-GMM (\ar{N} états, \ar{K} gauss.) & \ar{\%} & \ar{ } & \ar{ms} & \ar{ } \\
\bottomrule
\end{tabularx}
\caption{Résultats sur les 1\,000 fichiers de test (scripts \texttt{exp01}, \texttt{exp07}, \texttt{exp08}).}\label{tab:principal}
\end{table}

\begin{aremplir}
Commentaire (5 à 8 lignes) : quelle méthode est la meilleure, de combien, et est-ce significatif (intervalles de confiance, ou test de McNemar si l'écart est faible) ? Que montre l'écart entre corrélation et DTW sur l'intérêt de l'alignement non linéaire ? Que montre l'écart entre DTW et HMM sur l'intérêt d'un modèle statistique ?
\end{aremplir}

\figureres{exp04_caracteristiques.png}{0.65}{Effet de la représentation sur l'accuracy de la DTW (script \texttt{exp04\_caracteristiques.m}).}

\begin{aremplir}
Commentaire sur l'apport de la DCT (FBANK contre MFCC), des deltas et de la CMVN, et sur le choix Hann contre Hamming.
\end{aremplir}

\subsection{Mots confondus}

\begin{figure}[H]\centering
\fcolorbox{bordjaune}{jaune}{\parbox{0.8\linewidth}{\centering\small\vspace{1.5em}\textbf{FIGURE À INSÉRER} : matrice de confusion de la meilleure méthode (\texttt{evaluer(..., true)}, puis enregistrer la figure)\vspace{1.5em}}}
\caption{Matrice de confusion de la meilleure méthode classique.}
\end{figure}

\begin{aremplir}
Les paires les plus confondues (on s'attend par exemple à « no » et « go », « on » et « off », « up » et « off ») et leur explication phonétique, en s'appuyant sur les spectrogrammes de \texttt{figuresSignal.m}. Précision, rappel et F1 des mots les moins bien reconnus.
\end{aremplir}

\subsection{Robustesse au bruit}

\figureres{exp02_snr.png}{0.7}{Accuracy de la DTW en fonction du SNR, pour un bruit blanc et un bruit de robinet (script \texttt{exp02\_robustesse\_snr.m}).}

\begin{aremplir}
À partir de quel SNR l'accuracy chute-t-elle nettement ? Quel bruit est le plus gênant et pourquoi (répartition de son énergie en fréquence) ? Si le temps le permet, même courbe pour le HMM-GMM, et comparaison des deux méthodes.
\end{aremplir}

\subsection{Robustesse au locuteur}

\figureres{exp03_locuteur.png}{0.5}{Accuracy sur des locuteurs déjà entendus (dépendant) et sur des locuteurs nouveaux (indépendant) (script \texttt{exp03\_dependance\_locuteur.m}).}

\begin{aremplir}
Écart entre les deux tests pour la DTW (et pour les HMM si mesuré). Interprétation : une méthode qui mémorise des exemples (DTW) est-elle plus sensible au locuteur qu'une méthode qui apprend un modèle (HMM) ?
\end{aremplir}

\subsection{Coût de calcul}

\figureres{exp05_references.png}{0.8}{Accuracy et temps de reconnaissance de la DTW selon le nombre de références par mot (script \texttt{exp05\_nombre\_references.m}).}

\begin{aremplir}
Compromis précision et coût de la DTW. Résultats de \texttt{exp06\_cout.m} : temps d'extraction, temps de reconnaissance, facteur temps réel, mémoire des références (en Ko). Comparaison avec le coût des HMM, qui ne dépend pas du nombre d'exemples d'apprentissage.
\end{aremplir}

\subsection{Pires cas}

\begin{aremplir}
Quelques exemples d'erreurs écoutés et expliqués : VAD qui coupe le mot, locuteur atypique, enregistrement très bruité ou tronqué. Ces cas illustrent les limites des hypothèses de la section~\ref{sec:methodes}.
\end{aremplir}

%=====================================================================
\section{Benchmark des méthodes modernes}\label{sec:benchmark}\redaction{C}

Conformément à l'énoncé, les méthodes modernes ne sont pas réimplémentées : nous comparons nos résultats aux résultats publiés. Les chiffres ci-dessous proviennent des tableaux des articles originaux ; ils ne sont comparables qu'à l'intérieur d'une même tâche et d'un même protocole.

\begin{table}[H]\centering\small
\begin{tabularx}{\linewidth}{>{\raggedright\arraybackslash}p{4.6cm}cc>{\raggedright\arraybackslash}p{2.6cm}L}
\toprule
Méthode (référence) & Version & Accuracy & Taille & Remarque \\
\midrule
CNN de référence (Warden, 2018) & v2 & 88,2 \% & non précisée & baseline officielle \\
DS-CNN S (Zhang et al., 2017) & v1 & 94,4 \% & 38,6 Ko de mémoire & quantifié sur 8 bits \\
DS-CNN L (Zhang et al., 2017) & v1 & 95,4 \% & 497,6 Ko de mémoire & quantifié sur 8 bits \\
BC-ResNet-1 (Kim et al., 2021) & v2 & 96,9 \% & 9,2 k paramètres & CNN très compact \\
BC-ResNet-8 (Kim et al., 2021) & v2 & 98,7 \% & 321 k paramètres & état de l'art CNN \\
KWT-3 (Berg et al., 2021) & v2 & 98,56 \% & 5,4 M paramètres & Transformer, distillation \\
wav2vec 2.0 BASE, gelé (SUPERB) & v1 & 96,23 \% & 95 M paramètres & représentations auto-supervisées \\
HuBERT BASE, gelé (SUPERB) & v1 & 96,30 \% & 95 M paramètres & représentations auto-supervisées \\
\midrule
Notre meilleure méthode classique & v2, 10 mots & \ar{\%} & \ar{ } & protocole différent (10 classes, sans silence ni unknown) \\
\bottomrule
\end{tabularx}
\caption{Speech Commands, protocole à 12 classes (10 mots, silence, inconnu), résultats publiés.}\label{tab:kws}
\end{table}

\begin{table}[H]\centering\small
\begin{tabularx}{\linewidth}{>{\raggedright\arraybackslash}p{5cm}ccL}
\toprule
Méthode (référence) & Paramètres & WER test-clean / test-other & Conditions \\
\midrule
Conformer S (Gulati et al., 2020) & 10,3 M & 2,7 / 6,3 \% & sans modèle de langue \\
Conformer L (Gulati et al., 2020) & 118,8 M & 1,9 / 3,9 \% & avec modèle de langue \\
wav2vec 2.0 LARGE (Baevski et al., 2020) & 317 M & 1,8 / 3,3 \% & 960 h transcrites + modèle de langue \\
wav2vec 2.0 LARGE, 10 min transcrites & 317 M & 4,8 / 8,2 \% & 53\,000 h non transcrites \\
HuBERT X-LARGE (Hsu et al., 2021) & 964 M & 1,8 / 2,9 \% & 960 h transcrites + modèle de langue \\
Whisper large-v2 (Radford et al., 2023) & 1\,550 M & 2,7 / 5,2 \% & sans aucun entraînement sur LibriSpeech \\
\bottomrule
\end{tabularx}
\caption{Parole continue, LibriSpeech (phrases lues en anglais).}\label{tab:libri}
\end{table}

Le passage des GMM aux réseaux profonds a marqué une rupture : sur Switchboard (conversations téléphoniques, 309 h), le WER est passé de 27,4 \% pour le meilleur GMM-HMM à 18,5 \% pour le DNN-HMM, soit une réduction relative de 33 \% (Hinton et al., 2012) ; les TDNN de Kaldi ont apporté une amélioration relative moyenne supplémentaire de 5,5 \% (Peddinti et al., 2015).

\begin{aremplir}
Analyse critique (une demi-page) : position de vos méthodes par rapport aux résultats publiés, en rappelant les différences de protocole ; compromis performance / nombre de paramètres / données nécessaires (par exemple, BC-ResNet-1 atteint 96,9 \% avec 9\,200 paramètres, à comparer avec la mémoire de votre DTW) ; ce que les méthodes modernes apportent face au bruit et au locuteur ; ce que les méthodes classiques conservent (interprétabilité, peu de données, simplicité). Si l'exemple MathWorks de réseau convolutif a été lancé sur votre test, ajouter son résultat au tableau~\ref{tab:principal} : c'est la seule comparaison à protocole identique.
\end{aremplir}

%=====================================================================
\section{Perspectives d'applications (axe 3)}\label{sec:perspectives}\redaction{C}

\begin{aremplir}
Au moins deux perspectives (1 à 2 pages au total), choisies parmi celles du document 06 (section 7) ou proposées par vous. Pour chacune, un paragraphe structuré :
\begin{itemize}[leftmargin=*,itemsep=0pt]
\item domaine visé et besoin ;
\item contraintes d'exploitation (temps réel, embarqué, énergie, réglementation, confidentialité) ;
\item méthode recommandée et pourquoi, \textbf{en citant au moins un chiffre de vos expériences} ;
\item niveau de maturité de la méthode ;
\item verrous restant à lever.
\end{itemize}
Pistes proposées : commandes vocales dans les langues locales du Sénégal (peu de données : DTW avec peu de références, ou auto-supervisé) ; assistance aux personnes à mobilité réduite (fiabilité, dépendant du locuteur) ; commande mains libres en milieu industriel bruité (courbes SNR) ; mot de réveil sur objet connecté (mémoire, TinyML).
\end{aremplir}

%=====================================================================
\section{Conclusion}\redaction{tous}

\begin{aremplir}
Réponse explicite à la question de recherche, appuyée sur vos chiffres : quelle méthode classique résiste le mieux au bruit et à la variabilité du locuteur, dans quelles limites, et à quel coût ? Comment se situe-t-elle face aux méthodes modernes ? Limites de l'étude (sous-ensemble de 10 mots, conditions non appariées, PLP non implémentés...) et suites possibles.
\end{aremplir}

%=====================================================================
\begin{thebibliography}{99}\small
\bibitem{warden} P. Warden, « Speech Commands: A Dataset for Limited-Vocabulary Speech Recognition », arXiv:1804.03209, 2018.
\bibitem{davis} S. Davis, P. Mermelstein, « Comparison of parametric representations for monosyllabic word recognition in continuously spoken sentences », \emph{IEEE Transactions on Acoustics, Speech, and Signal Processing}, 1980.
\bibitem{hermansky} H. Hermansky, « Perceptual linear predictive (PLP) analysis of speech », \emph{Journal of the Acoustical Society of America}, 1990.
\bibitem{sakoe} H. Sakoe, S. Chiba, « Dynamic programming algorithm optimization for spoken word recognition », \emph{IEEE Transactions on Acoustics, Speech, and Signal Processing}, 1978.
\bibitem{rabiner} L. Rabiner, « A tutorial on hidden Markov models and selected applications in speech recognition », \emph{Proceedings of the IEEE}, 1989.
\bibitem{rabsam} L. Rabiner, M. Sambur, « An algorithm for determining the endpoints of isolated utterances », \emph{Bell System Technical Journal}, 1975.
\bibitem{hinton} G. Hinton et al., « Deep Neural Networks for Acoustic Modeling in Speech Recognition: The Shared Views of Four Research Groups », \emph{IEEE Signal Processing Magazine}, 2012.
\bibitem{peddinti} V. Peddinti, D. Povey, S. Khudanpur, « A time delay neural network architecture for efficient modeling of long temporal contexts », Interspeech, 2015.
\bibitem{gulati} A. Gulati et al., « Conformer: Convolution-augmented Transformer for Speech Recognition », Interspeech, 2020.
\bibitem{baevski} A. Baevski, H. Zhou, A. Mohamed, M. Auli, « wav2vec 2.0: A Framework for Self-Supervised Learning of Speech Representations », NeurIPS, 2020.
\bibitem{hsu} W.-N. Hsu et al., « HuBERT: Self-Supervised Speech Representation Learning by Masked Prediction of Hidden Units », \emph{IEEE/ACM Transactions on Audio, Speech, and Language Processing}, 2021.
\bibitem{yang} S. Yang et al., « SUPERB: Speech processing Universal PERformance Benchmark », Interspeech, 2021.
\bibitem{radford} A. Radford et al., « Robust Speech Recognition via Large-Scale Weak Supervision », ICML, 2023.
\bibitem{zhang} Y. Zhang, N. Suda, L. Lai, V. Chandra, « Hello Edge: Keyword Spotting on Microcontrollers », arXiv:1711.07128, 2017.
\bibitem{kim} B. Kim, S. Chang, J. Lee, D. Sung, « Broadcasted Residual Learning for Efficient Keyword Spotting », Interspeech, 2021.
\bibitem{berg} A. Berg, M. O'Connor, M. Tairum Cruz, « Keyword Transformer: A Self-Attention Model for Keyword Spotting », Interspeech, 2021.
\end{thebibliography}

\end{document}
